% ---------
% --------
%!TEX TS-program = pdflatex
%!TEX encoding = UTF-8 Unicode
%!TEX spellcheck = English (Aspell)

\documentclass[11pt]{article}
\usepackage{jeffe,handout,graphicx}
\usepackage{fourier-orns}
\usepackage{mathtools}
\usepackage[normalem]{ulem} % either use this (simple) or
\usepackage{tikz}

\def\To{\leadsto}
\def\Tostar{\mathrel{\To\!\!\!^*}}
\def\derives{\Tostar}

\def\Cdot{\mathbin{\text{\normalfont \textbullet}}}
\def\Sym#1{\texttt{\color{BrickRed}#1}}
\def\BlueSym#1{\textbf{\texttt{\color{RoyalBlue}#1}}}
\allowdisplaybreaks

\begin{document}

\begin{center}
\Large\textbf{CSCI 432/532, Spring 2026}%
\\
\LARGE\textbf{Homework 6}%
\\[0.5ex]
\large Due Monday, March 2  at 8:30am
\end{center}

\bigskip
\hrule
\bigskip

\subsection*{Submission Requirements}
\begin{itemize}
    \item Type or clearly hand-write your solutions into a PDF format so that they are legible and professional. Submit your PDF on Canvas. 
	\item Put each \textbf{sub}problem on a separate page and select the corresponding problem via the Gradescope interface.
    \item Do not submit your first draft. Type or clearly re-write your solutions for your final submission. If your submission is not legible, we will ask you to resubmit.
\end{itemize}

\subsection*{Collaboration and Resources}

The \emph{best} resources for completing the homework are (in no particular order):
\begin{itemize}
	\item Your own notes from lectures and problem sessions, and/or lecture recordings.
	\item The lecture notes linked from the course website for the lecture day.
	\item The questions channel on Discord.
	\item Office hours.
	\item Discussions with classmates.
\end{itemize}
You may also access almost any resource in preparing your solution to the
homework. The exception is that you may not seek answers to the problems on the
homework directly, such as by pasting them into a GenAI tool, searching for
solutions on a search engine, looking for them on Chegg or similar websites, or
asking another person for the solution.

Additionally, you \EMPH{must}
\begin{itemize}
    \item Write your solutions in your own words---this means that you should never be \emph{copying}
    anything of substance from any source, and
    \item credit every resource you use, including GenAI tools, by providing links or full chat transcripts.
    If you use the provided LaTeX template, you can use the \texttt{Sources}
    environment for this. Ask if you need help!
\end{itemize}


\begin{boxquote}{0.9}
Remember, if you choose to use GenAI tools, not only must you provide a full
transcript of your conversation (or a link to one), you must also use the
following prompt (or some variation of it), and provide a full transcript of
the conversation.
\end{boxquote}  

Prompt:
\begin{boxquote}{0.9}
You are acting as a teaching assistant for an advanced undergraduate/graduate
algorithms and computer science theory course.

Your role is not to provide full solutions or final answers.

Instead, you should act like we are having a conversation. Ask me a single
question and let me respond. Avoid asking me many questions at once or giving
me a lot of information at one time. Guide me using hints, leading questions,
partial insights, and sanity checks.  Help me debug my own proofs, algorithms,
or reductions. Do not give a complete solution or full proofs. Be patient with
me. Don't give me details so that I can move on to the next part of a problem.
Make sure that I understood before moving on.

The course uses Jeff Erickson's Models of Computation lecture notes, so you should
guide me to answer problems using the definitions, notation, and style from those notes.

Assume I am a serious student trying to learn, not to shortcut the assignment.
I am attaching the assignment, so you should refuse to provide a complete solution
to the problems listed. Of course, you can walk me through the "Solved
Problems" listed at the end if I ask.
\end{boxquote}  



\newpage
%----------------------------------------------------------------------

\headers{CSCI 432/532}{Homework 6 (due March 2)}{Spring 2026}

\begin{enumerate}
%\parindent 1.5em \itemsep 4ex plus 0.5fil

%----------------------------------------------------------------------
\item Describe (as in the problem session solutions---note that you do
\emph{not} need to give a full specification of states, transitions, etc.) a
Turing machine that decides the following language: $\{ww : w \in \{\Sym0,
\Sym1\}^*\}$.

\begin{Rubric}{TM rubric}
10 points =
\begin{itemize}
\item[+ 2] for an English description of a Turing machine, regardless of correctness.
\item[+ 6] for correctness.
\item[+ 2] for an appropriate mix of detail and high-level intuition.
\end{itemize}
\end{Rubric}

%----------------------------------------------------------------------
\item A \EMPH{Turing machine with stay put instead of left} is similar to an
ordinary Turing machine, but the transition function has the form $\delta: Q
\times \Gamma \to Q \times \Gamma \times \{R, S\}$. At each point, the machine
can move its head right or let it stay put in the same position.
\begin{enumerate}[(a)]
	\item Show that this Turing machine variant is \emph{not}
equivalent to a standard Turing machine. 
\item What class of languages do these
machines recognize? 
\end{enumerate}

\begin{Rubric}{TM simulation rubric}
10 points =
\begin{itemize}
\item[+ 2] for an English description of either (1) how to simulate a standard Turing machine with the new type of Turing machine or (2) why standard Turing machine can't be simulated with the new type of Turing machine.
\item[+ 6] for correctness.
\item[+ 2] for an appropriate mix of detail and high-level intuition.
\end{itemize}
\end{Rubric}

\item 
Consider the language $\mathsc{NeverReject} = \{ \langle M \rangle :
\textsc{Reject}(M) = \emptyset \}$; that is, the set of all (encodings of)
Turing machines that do not reject any input. 
\begin{enumerate}
	\item Prove that $\mathsc{NeverReject}$
is undecidable via \EMPH{diagonalization}.
\begin{Rubric}{Diagnoalization rubric}
10 points =
\begin{itemize}
\item[+ 4] for a correct wrapper Turing machine
\item[+ 6] for self-contradiction proof (= 3 for $\Longleftarrow$ + 3 for $\Longrightarrow$)
\end{itemize}
\end{Rubric}

\item Prove that $\mathsc{NeverReject}$ is undecidable via \EMPH{reduction}
(for example, from the halting problem).

\begin{Rubric}{Reduction rubric}
10 points =
\begin{itemize}
\item[+ 4] for a correct reduction
\item[+ 3] for ''if" proof
\item[+ 3] for ''only if" proof
\end{itemize}
\end{Rubric}


\end{enumerate}
\end{enumerate}


% ===========================================================================
\newpage

\subsection*{Solved problems}

\begin{enumerate}\parindent 1.5em

\item A \EMPH{Turing machine with a doubly infinite tape} is similar to an ordinary
Turing machine, but the tape is infinite in both directions. Show that this Turing
machine variant is equivalent to a standard Turing machine.

\begin{Solution}
To show that a Turing machine with a doubly infinite tape is equivalent
to a standard Turing machine, we can show that a standard Turing
machine can simulate a Turing machine with a doubly infinite tape.
To do so, we must simulate being able to write arbitrarily far to the left of the starting
point. In class, we saw that this can be done by simulating a second tape representing
negative indices using our tape alphabet.

Another way to do this is to use the original tape alphabet, but copy the entire
tape contents one position to the right whenever we need to write to the left of
the starting point. We can add a left sentinel symbol and right sentinel symbol to the
tape alphabet to mark the left and right ends of the tape contents. Whenever we need
to write to the right sentinel, we shift it one to the right and write the new symbol
in its old place. Whenever we need to write to the left sentinel, we shift all the tape contents one
to the right and write the new symbol in the left sentinel's old place. This way, we can simulate a doubly infinite tape using a standard Turing machine.
\end{Solution}

\item Conside the language $\mathsc{SelfHalt} = \{ \langle M \rangle : M \text{
halts on input } \langle M \rangle \}$; that is, the set of all (encodings of)
Turing machines that halt when given their own encoding as input. Prove that
$\mathsc{SelfHalt}$ is undecidable.

\begin{Solution}
Suppose to the contrary that there is a Turing machine $SH$ such that
$\mathsc{Accept}(SH) = \mathsc{SelfHalt}$ and $\mathsc{Diverge}(SH) =
\emptyset$.

Let $SH^X$ be the Turing machine
obtained from $SH$ by redirecting every transition to \Sym{accept} to a new hanging
state \Sym{hang}, and then redirecting every transition to \Sym{reject} to
\Sym{accept}. Then $\mathsc{Accept}(SH^X) = \Sigma^* \setminus \mathsc{SelfHalt}$ and
$\mathsc{Reject}(SH^X) = \emptyset$.  It follows that $SH^X$ accepts $\langle SH^X
\rangle$ if and only if $SH^X$ does not halt on $\langle SH^X \rangle$, and
we have a contradiction.
\end{Solution}

\item
Consider the language $\mathsc{SometimesHalt} = \set{ \seq{M} \mid \text{$M$~halts
on at least one input string}}$.  Note that $\seq{M}\in \mathsc{SometimesHalt}$ does not imply that $M$ \emph{accepts} any strings; it is enough that $M$ \emph{halts} on (and possibly rejects) some string.

Prove that \mathsc{SometimesHalt} is undecidable.

\begin{Solution}
We can reduce the standard halting problem to \mathsc{Sometimes\-Halt} as follows:
\begin{algo}
	\textul{$\mathsc{DecideHalt}(\seq{M, w})$:}\+
\\	Encode the following Turing machine $M'$:\+
\\	\color{OliveGreen}
	\qquad \begin{algorithm}
		\textul{$M'(x)$:}\+
	\\		(ignore $x$)
	\\		run $M$ on input $w$
	\end{algorithm}
\\[1ex]
	return $\mathsc{DecideSometimesHalt}(\seq{M'})$\-
\end{algo}
We prove this reduction correct as follows:

\begin{itemize}
\item[$\Longrightarrow$]
Suppose $M$ halts on input $w$.\\
Then $M'$ halts on \emph{every} input string $x$.\\
So \mathsc{DecideSometimesHalt} must accept the encoding $\seq{M’}$.\\
We conclude that \mathsc{DecideHalt} correctly accepts the encoding $\seq{M, w}$.

\item[$\Longleftarrow$]
Suppose $M$ does not halt on input $w$.\\
Then $M'$ diverges on \emph{every} input string $x$.\\
So \mathsc{DecideSometimesHalt} must reject the encoding $\seq{M’}$.\\
We conclude that \mathsc{DecideHalt} correctly rejects the encoding $\seq{M, w}$.
\end{itemize} 
\end{Solution}

\end{enumerate}
\end{document}

